% Impact Analysis of Intrusion Detection Systems on System Performance
% !TeX spellcheck = en-US
% LTeX: language=en-US
% !TeX encoding = utf8
% !TeX program = lualatex
% !BIB program = bibtex
% -*- coding:utf-8 mod:LaTeX -*-

\documentclass[runningheads,a4paper,english]{llncs}[2022/01/12]

\usepackage{iftex}

% backticks (`) are rendered as such in verbatim environments.
% See following links for details:
%   - https://tex.stackexchange.com/a/341057/9075
%   - https://tex.stackexchange.com/a/47451/9075
%   - https://tex.stackexchange.com/a/166791/9075
\usepackage{upquote}

% Set English as language and allow to write hyphenated"=words
%
% Even though `american`, `english` and `USenglish` are synonyms for babel package (according to https://tex.stackexchange.com/questions/12775/babel-english-american-usenglish), the llncs document class is prepared to avoid the overriding of certain names (such as "Abstract." -> "Abstract" or "Fig." -> "Figure") when using `english`, but not when using the other 2.
% english has to go last to set it as default language
\usepackage[ngerman,main=english]{babel}
%
% Hint by http://tex.stackexchange.com/a/321066/9075 -> enable "= as dashes
\addto\extrasenglish{\languageshorthands{ngerman}\useshorthands{"}}

% Links behave as they should. Enables "\url{...}" for URL typesettings.
% Allow URL breaks also at a hyphen, even though it might be confusing: Is the "-" part of the address or just a hyphen?
% See https://tex.stackexchange.com/a/3034/9075.
\usepackage[hyphens]{url}

% When activated, use text font as url font, not the monospaced one.
% For all options see https://tex.stackexchange.com/a/261435/9075.
% \urlstyle{same}

% Improve wrapping of URLs - hint by http://tex.stackexchange.com/a/10419/9075
\makeatletter
\g@addto@macro{\UrlBreaks}{\UrlOrds}
\makeatother

% nicer // - solution by http://tex.stackexchange.com/a/98470/9075
% DO NOT ACTIVATE -> prevents line breaks
%\makeatletter
%\def\Url@twoslashes{\mathchar`\/\@ifnextchar/{\kern-.2em}{}}
%\g@addto@macro\UrlSpecials{\do\/{\Url@twoslashes}}
%\makeatother

%% !!! If you change the font, be sure that words such as "workflow" can
%% !!! still be copied from the PDF. If this is not the case, you have
%% !!! to use glyphtounicode. See comment at cmap package.
%%
%% Background: "workflow" contains "fl" which is a ligature, which in turn
%%             is rendered as one character in the PDF and needs to be split
%%             whily copying.

\ifluatex
  \usepackage[no-math]{fontspec}
  \usepackage{unicode-math}

  % Typewriter font (for source code etc)
  % Use New Computer Modern font (Computer Modern is the default LaTeX font; this is the implemented modern variant)
  % Source: https://tug.org/FontCatalogue/newcomputermoderntypewriter/

  \setmainfont[
    ItalicFont=NewCM10-Italic.otf,
    BoldFont=NewCM10-Bold.otf,
    BoldItalicFont=NewCM10-BoldItalic.otf,
    SmallCapsFeatures={Numbers=OldStyle}]{NewCM10-Regular.otf}

  \setsansfont[
    ItalicFont=NewCMSans10-Oblique.otf,
    BoldFont=NewCMSans10-Bold.otf,
    BoldItalicFont=NewCMSans10-BoldOblique.otf,
    SmallCapsFeatures={Numbers=OldStyle}]{NewCMSans10-Regular.otf}

  \setmonofont[ItalicFont=NewCMMono10-Italic.otf,
    BoldFont=NewCMMono10-Bold.otf,
    BoldItalicFont=NewCMMono10-BoldOblique.otf,
    SmallCapsFeatures={Numbers=OldStyle}]{NewCMMono10-Regular.otf}

  \setmathfont{NewCMMath-Regular.otf}

  % Enable proper ligatures
  % For more information see https://ctan.org/pkg/selnolig
  % language "english" or "ngerman" is passed to selnolig by the document class
  \usepackage{selnolig}

\else
  % This is the modern package for "Computer Modern".
  % In case this gets activated, one has to switch from cmap package to glyphtounicode (in the case of pdflatex)
  \usepackage[%
    rm={oldstyle=false,proportional=true},%
    sf={oldstyle=false,proportional=true},%
    % By using 'variable=true' the monospaced font can be used as variable font (with differents widths per letter)
    % However, this makes listings look ugly.
    tt={oldstyle=false,proportional=true,variable=false},%
    qt=false%
  ]{cfr-lm}

  % Has to be loaded AFTER any font packages. See https://tex.stackexchange.com/a/2869/9075.
  \usepackage[T1]{fontenc}
\fi

% Character protrusion and font expansion. See http://www.ctan.org/tex-archive/macros/latex/contrib/microtype/

\usepackage[
  babel=true, % Enable language-specific kerning. Take language-settings from the languge of the current document (see Section 6 of microtype.pdf)
  expansion=alltext,
  protrusion=alltext-nott, % Ensure that at listings, there is no change at the margin of the listing
  % In the standard configuration, this template is always in the final mode, so this option only makes a difference if "pros" use the draft mode
  final % Always enable microtype, even if in draft mode. This helps finding bad boxes quickly.
]{microtype}

% \texttt{test -- test} keeps the "--" as "--" (and does not convert it to an en dash)
\DisableLigatures{encoding = T1, family = tt* }

%\DeclareMicrotypeSet*[tracking]{my}{ font = */*/*/sc/* }%
%\SetTracking{ encoding = *, shape = sc }{ 45 }
% Source: http://homepage.ruhr-uni-bochum.de/Georg.Verweyen/pakete.html
% Deactiviated, because does not look good

\usepackage{graphicx}

% Diagonal lines in a table - http://tex.stackexchange.com/questions/17745/diagonal-lines-in-table-cell
% Slashbox is not available in texlive (due to licensing) and also gives bad results. Thus, we use diagbox
\usepackage{diagbox}

\ifluatex
  \usepackage{spelling}
  \spellingoutput{off}
\fi

\usepackage[dvipsnames, table]{xcolor}
% Code Listings
\usepackage{listings}

\definecolor{eclipseStrings}{RGB}{42,0.0,255}
\definecolor{eclipseKeywords}{RGB}{127,0,85}
\colorlet{numb}{magenta!60!black}

% JSON definition
% Source: https://tex.stackexchange.com/a/433961/9075

\lstdefinelanguage{json}{
  basicstyle=\normalfont\ttfamily,
  commentstyle=\color{eclipseStrings}, % style of comment
  stringstyle=\color{eclipseKeywords}, % style of strings
  numbers=left,
  numberstyle=\scriptsize,
  stepnumber=1,
  numbersep=8pt,
  showstringspaces=false,
  breaklines=true,
  frame=lines,
  % backgroundcolor=\color{gray}, %only if you like
  string=[s]{"}{"},
  comment=[l]{:\ "},
  morecomment=[l]{:"},
  literate=
    *{0}{{{\color{numb}0}}}{1}
    {1}{{{\color{numb}1}}}{1}
    {2}{{{\color{numb}2}}}{1}
    {3}{{{\color{numb}3}}}{1}
    {4}{{{\color{numb}4}}}{1}
    {5}{{{\color{numb}5}}}{1}
    {6}{{{\color{numb}6}}}{1}
    {7}{{{\color{numb}7}}}{1}
    {8}{{{\color{numb}8}}}{1}
    {9}{{{\color{numb}9}}}{1}
}

\lstset{
  % everything between (* *) is a latex command
  escapeinside={(*}{*)},
  %
  language=json,
  %
  showstringspaces=false,
  %
  basicstyle=\footnotesize\ttfamily,
  %
  commentstyle=\slshape,
  %
  % default: \rmfamily
  stringstyle=\ttfamily,
  %
  breaklines=true,            % Zeilen werden umbrochen
  %
  breakatwhitespace=true,
  %
  % alternative: fixed
  columns=flexible,
  %
  tabsize=2,                  % Groesse von Tabs
  %
  numbers=left,
  %
  numberstyle=\tiny,
  %
  basewidth=.5em,
  %
  xleftmargin=.5cm,
  %
  % aboveskip=0mm,
  %
  % belowskip=0mm,
  %
  captionpos=b
}

\ifpdftex
  % Enable Umlauts when using \lstinputputlisting.
  % See https://stackoverflow.com/a/29260603/873282 and https://tex.stackexchange.com/a/24532/9075 for details.
  % listingsutf8 did not work in June 2020.
  \lstset{extendedchars=true, literate=
    {á}{{\'a}}1 {é}{{\'e}}1 {í}{{\'i}}1 {ó}{{\'o}}1 {ú}{{\'u}}1
  {Á}{{\'A}}1 {É}{{\'E}}1 {Í}{{\'I}}1 {Ó}{{\'O}}1 {Ú}{{\'U}}1
  {à}{{\`a}}1 {è}{{\`e}}1 {ì}{{\`i}}1 {ò}{{\`o}}1 {ù}{{\`u}}1
  {À}{{\`A}}1 {È}{{\'E}}1 {Ì}{{\`I}}1 {Ò}{{\`O}}1 {Ù}{{\`U}}1
  {ä}{{\"a}}1 {ë}{{\"e}}1 {ï}{{\"i}}1 {ö}{{\"o}}1 {ü}{{\"u}}1
  {Ä}{{\"A}}1 {Ë}{{\"E}}1 {Ï}{{\"I}}1 {Ö}{{\"O}}1 {Ü}{{\"U}}1
  {â}{{\^a}}1 {ê}{{\^e}}1 {î}{{\^i}}1 {ô}{{\^o}}1 {û}{{\^u}}1
  {Â}{{\^A}}1 {Ê}{{\^E}}1 {Î}{{\^I}}1 {Ô}{{\^O}}1 {Û}{{\^U}}1
  {Ã}{{\~A}}1 {ã}{{\~a}}1 {Õ}{{\~O}}1 {õ}{{\~o}}1
  {œ}{{\oe}}1 {Œ}{{\OE}}1 {æ}{{\ae}}1 {Æ}{{\AE}}1 {ß}{{\ss}}1
  {ű}{{\H{u}}}1 {Ű}{{\H{U}}}1 {ő}{{\H{o}}}1 {Ő}{{\H{O}}}1
  {ç}{{\c c}}1 {Ç}{{\c C}}1 {ø}{{\o}}1 {å}{{\r a}}1 {Å}{{\r A}}1
  }
\fi

\lstloadlanguages{% Check dokumentation for further languages...
  %[Visual]Basic
  %Pascal
  %C
  %C++
  %XML
  %HTML
}

% For easy quotations: \enquote{text}
% This package is very smart when nesting is applied, otherwise textcmds (see below) provides a shorter command
\usepackage[autostyle=true]{csquotes}

% Enable using "`quote"' - see https://tex.stackexchange.com/a/150954/9075
\defineshorthand{"`}{\openautoquote}
\defineshorthand{"'}{\closeautoquote}

% Nicer tables (\toprule, \midrule, \bottomrule)
\usepackage{booktabs}

% Extended enumerate, such as \begin{compactenum}
\usepackage{paralist}

% Bibliopgraphy enhancements
%  - enable \cite[prenote][]{ref}
%  - enable \cite{ref1,ref2}
% Alternative: \usepackage{cite}, which enables \cite{ref1, ref2} only (otherwise: Error message: "White space in argument")

% Doc: http://texdoc.net/natbib
\usepackage[%
  square,        % for square brackets
  comma,         % use commas as separators
  numbers,       % for numerical citations;
  %sort           % orders multiple citations into the sequence in which they appear in the list of references;
  sort&compress  % as sort but in addition multiple numerical citations are compressed if possible (as 3-6, 15);
]{natbib}

% In the bibliography, references have to be formatted as 1., 2., ... not [1], [2], ...
\renewcommand{\bibnumfmt}[1]{#1.}

% Enable hyperlinked author names in the case of \citet
% Source: https://tex.stackexchange.com/a/76075/9075
\usepackage{etoolbox}
\makeatletter
\patchcmd{\NAT@test}{\else \NAT@nm}{\else \NAT@hyper@{\NAT@nm}}{}{}
\makeatother

% Prepare more space-saving rendering of the bibliography
% Source: https://tex.stackexchange.com/a/280936/9075
\SetExpansion
[ context = sloppy,
  stretch = 30,
  shrink = 60,
  step = 5 ]
{ encoding = {OT1,T1,TS1} }
{ }

% Put figures aside a text
% Even though the package is from 1998, it works well
\usepackage[rflt]{floatflt}

% Farbige Tabellen
% ----------------
% Das Paket colortbl wird inzwischen automatisch durch xcolor geladen
%
% Erweiterte Funktionen innerhalb von Tabellen
% --------------------------------------------
%%% Doc: http://mirror.ctan.org/tex-archive/macros/latex/contrib/multirow/multirow.sty
\usepackage{multirow} % Mehrfachspalten
%
%%% Doc: Documentation inside dtx Package
\usepackage{dcolumn}  % Ausrichtung an Komma oder Punkt

%%% Doc: http://mirror.ctan.org/tex-archive/macros/latex/contrib/supertabular/supertabular.pdf
%\usepackage{supertabular}

%%% Fussnoten/Endnoten ===================================================

% EN: Put footnotes below floats
% DE: Fußnoten unter Gleitumgebungen ("floats") platzieren
% Source: https://tex.stackexchange.com/a/32993/9075
\usepackage{stfloats}
\fnbelowfloat

% EN: Extended support for footnotes
% DE: Fußnoten
%
%\usepackage{dblfnote}  %Zweispaltige Fußnoten
%
% Keine hochgestellten Ziffern in der Fußnote (KOMA-Script-spezifisch):
%\deffootnote[1.5em]{0pt}{1em}{\makebox[1.5em][l]{\bfseries\thefootnotemark}}
%
% Abstand zwischen Fußnoten vergrößern:
%\setlength{\footnotesep}{.85\baselineskip}
%
% EN: Following command disables the separting line of the footnote
% DE: Folgendes Kommando deaktiviert die Trennlinie zur Fußnote
%\renewcommand{\footnoterule}{}
%
%\addtolength{\skip\footins}{\baselineskip} % Abstand Text <-> Fußnote

% DE: Fußnoten immer ganz unten auf einer \raggedbottom-Seite
% DE: fnpos kommt aus dem yafoot package
%\usepackage{fnpos}
%\makeFNbelow
%\makeFNbottom

% TODO (and comment) configuration
%
% - \todo (from todo, easy-todo, todonotes) / \TODO (from fixmetodonotes) - for "normal" TODOs
% - \todofix - "important" TODOs
%
% - \textcomment - highlights text and has a hover comment
% - \sidecomment - just puts a comment to the side. Note: \comment MUST NOT be used as command name, it is already defined by much packages (mathdesign, mindflow, verbatim, and others)
%
% - \missingfigure
%
% - \textmarker
% - \modified
% - \change      - adresses a review comment

% Enable nice comments
\usepackage{pdfcomment}

\newcommand{\textcomment}[2]{\colorbox{yellow!60}{#1}\pdfcomment[color={0.234 0.867 0.211},hoffset=-6pt,voffset=10pt,opacity=0.5]{#2}}

% Small PDF comment
% 1. Parameter: Comment
\newcommand{\sidecomment}[1]{\pdfcomment[color={0.045 0.278 0.643},voffset=4pt,icon=Note]{#1}}
% Disabled variant - for the final PDF
%\newcommand{\sidecomment}[1]{}

\newcommand{\todo}[1]{TODO!\sidecomment{#1}}

% Änderungen
%
% 1. Parameter: Review-Kommentar
% 2. Parameter: Neuer Text
\newcommand{\change}[2]{{\color{red}#2}\pdfcomment[color={0.234 0.867 0.211},voffset=8pt,opacity=0.5]{#1}}
% Disabled variant - for the final PDF
%\newcommand{\change}[2]{#2}

% Define default commands
\makeatletter
\@ifundefined{missingfigure}{\newcommand{\missingfigure}{... missing figure ...}}{}
\@ifundefined{textcomment}{\newcommand{\textcomment}[2]{#1 \todo{#2}}}{}
\@ifundefined{sidecomment}{\newcommand{\sidecomment}[1]{\marginpar{#1}}}{}
\@ifundefined{todo}{\newcommand{\todo}[1]{\sidecomment{#1}}}{}
\@ifundefined{TODO}{\newcommand{\TODO}[1]{\todo{#1}}}{}
\@ifundefined{todofix}{\newcommand{\todofix}[1]{\todo{#1}}}{}
\@ifundefined{change}{\newcommand{\change}[2]{#1 $\rightarrow$ #2}}{}
\makeatother

% Textmarker (Textfarbe rot)
\newcommand{\textmarker}[1]{{\color{red} #1}\xspace}

% Modified (Text blau)
\newcommand{\modified}[1]{{\color{blue!60!black} #1}\xspace}

\usepackage[group-minimum-digits=4,per-mode=fraction]{siunitx}

% Enable that parameters of \cref{}, \ref{}, \cite{}, ... are linked so that a reader can click on the number an jump to the target in the document
\usepackage{hyperref}

% Enable hyperref without colors and without bookmarks
\hypersetup{
  hidelinks,
  colorlinks=true,       % Links erhalten Farben statt Kaeten
  raiselinks=true,       % calculate real height of the link
  allcolors=black,
  pdfstartview=Fit,
  breaklinks=true,       % Links ueberstehen Zeilenumbruch
  hypertexnames=false,   % Fix jumping to algorithm line - http://tex.stackexchange.com/a/156404/9075
}

% Enable correct jumping to figures when referencing
\usepackage[all]{hypcap}

\usepackage[caption=false,font=footnotesize]{subfig}

% Alternative for making subfigures:
% Part of the caption package. See http://www.ctan.org/pkg/caption
% Ersetzt die Pakete subfigure und subfig - siehe https://tex.stackexchange.com/a/13778/9075
%
% (subfigure is outdated. subfig is maintained, but subcaption is better)
% See: http://tex.stackexchange.com/questions/13625/subcaption-vs-subfig-best-package-for-referencing-a-subfigure
%\usepackage[hypcap=true]{subcaption}

\usepackage{mindflow}

% Extensions for references inside the document (\cref{fig:sample}, ...)
% Enable usage \cref{...} and \Cref{...} instead of \ref: Type of reference included in the link
% That means, "Figure 5" is a full link instead of just "5".
\usepackage[capitalise,nameinlink]{cleveref}

\crefname{section}{Sect.}{Sect.}
\Crefname{section}{Section}{Sections}
\crefname{listing}{List.}{List.}
\crefname{listing}{Listing}{Listings}
\Crefname{listing}{Listing}{Listings}
\crefname{lstlisting}{Listing}{Listings}
\Crefname{lstlisting}{Listing}{Listings}

\usepackage{lipsum}

% For demonstration purposes only
% These packages can be removed when all examples have been deleted
\usepackage[math]{blindtext}
\usepackage{mwe}
\usepackage[realmainfile]{currfile}
\usepackage{tcolorbox}
\tcbuselibrary{listings}

%introduce \powerset - hint by http://matheplanet.com/matheplanet/nuke/html/viewtopic.php?topic=136492&post_id=997377
\DeclareFontFamily{U}{MnSymbolC}{}
\DeclareSymbolFont{MnSyC}{U}{MnSymbolC}{m}{n}
\DeclareFontShape{U}{MnSymbolC}{m}{n}{
  <-6>    MnSymbolC5
  <6-7>   MnSymbolC6
  <7-8>   MnSymbolC7
  <8-9>   MnSymbolC8
  <9-10>  MnSymbolC9
  <10-12> MnSymbolC10
  <12->   MnSymbolC12%
}{}
\DeclareMathSymbol{\powerset}{\mathord}{MnSyC}{180}

% Allows for defining commands that don't eat spaces.
\usepackage{xspace}
% Adds compatibility to \xspace und \enquote
\makeatletter
\xspaceaddexceptions{\grqq \grq \csq@qclose@i \} }
\makeatother

\newcommand{\eg}{e.g.,\ }
\newcommand{\ie}{i.e.,\ }

% Enable hyphenation at other places as the dash.
% Example: applicaiton\hydash specific
\makeatletter
\newcommand{\hydash}{\penalty\@M-\hskip\z@skip}
% Definition of "= taken from http://mirror.ctan.org/macros/latex/contrib/babel-contrib/german/ngermanb.dtx
\makeatother

% Add manual adapted hyphenation of English words
% See https://ctan.org/pkg/hyphenex and https://tex.stackexchange.com/a/22892/9075 for details
\input{ushyphex}

% correct bad hyphenation here
\hyphenation{
  op-tical net-works semi-conduc-tor
  % May not be hypphenated
  AROMA TOSCA BPMN OASIS OMG DMTF IT DevOps
}

\input{commands}

% Add copyright
%
% This is recommended if you intend to send the version to colleagues
% See https://ctan.org/pkg/llncsconf for details
\iffalse
  % state: intended | submitted | llncs
  % you can add "crop" if the paper should be cropped to the format Springer is publishing
  \usepackage[intended]{llncsconf}

  \conference{name of the conference}

  % in case of "proceedings" (final version!)
  % example: \llncs{Anonymous et al. (eds). \emph{Proceedings of the International Conference on \LaTeX-Hacks}, LNCS~42. Some Publisher, 2016.}{0042}
  % 0042 denotes an example start page
  \llncs{book editors and title}{0042}
\fi

\ifpdftex
  % Enable copy and paste of text from the PDF
  % Only required for pdflatex. It "just works" in the case of lualatex.
  % Alternative: cmap or mmap package
  % mmap enables mathematical symbols, but does not work with the newtx font set
  % See: https://tex.stackexchange.com/a/64457/9075
  % Other solutions outlined at http://goemonx.blogspot.de/2012/01/pdflatex-ligaturen-und-copynpaste.html and http://tex.stackexchange.com/questions/4397/make-ligatures-in-linux-libertine-copyable-and-searchable
  % Trouble shooting outlined at https://tex.stackexchange.com/a/100618/9075
  %
  % According to https://tex.stackexchange.com/q/451235/9075 this is the way to go
  \input{glyphtounicode}
  \pdfgentounicode=1
\fi
\begin{document}

\title{Performance Impact Analysis of Intrusion Detection \& Prevention Systems (IDS/IPS) on Windows 11: Study of Symantec Endpoint Security and OPNsense Firewall computational resource consumption}
\titlerunning{IDS/IPS Performance Impact on Windows 11}

% Single insitute
\author{Robert Molnar}

% If there are too many authors, use \authorrunning
%\authorrunning{First Author et al.}

\institute{Universitatea de Vest Timisoara}

%% Multiple insitutes - ALTERNATIVE to the above
% \author{%
%     Firstname Lastname\inst{1} \and
%     Firstname Lastname\inst{2}
% }
%
%If there are too many authors, use \authorrunning
%  \authorrunning{First Author et al.}
%
%  \institute{
%      Insitute 1\\
%      \email{...}\and
%      Insitute 2\\
%      \email{...}
%}

\maketitle

\begin{abstract}
This paper presents a comprehensive performance impact analysis of Intrusion Detection Systems (IDS) on Windows 11, specifically examining Symantec antivirus and OPNsense firewall. Through systematic benchmarking of four distinct system configurations (baseline without IDS, antivirus only, firewall only, and combined IDS), we quantify the computational overhead across six performance criteria: boot time, RAM consumption, process count, application launch performance, local network file transfers via SMB, and remote file downloads via FTP. Results from 5 iterations per test demonstrate measurable performance degradation with IDS deployment, with combined antivirus and firewall showing the highest impact. This study provides empirical evidence for system administrators and security professionals to make informed decisions about IDS deployment strategies while balancing security and performance requirements.

\keywords{Intrusion Detection Systems \and Performance Benchmarking \and Symantec \and OPNsense \and Windows 11 \and System Performance}
\end{abstract}


\section{Introduction}
\label{sec:introduction}

Intrusion Detection Systems (IDS), including antivirus software and firewalls, are essential security components in modern computing environments. However, these protective mechanisms impose computational overhead that can impact system performance. Understanding this trade-off between security and performance is critical for system administrators, IT professionals, and end users who must balance protection requirements with operational efficiency.

While antivirus and firewall software vendors often claim minimal performance impact, independent empirical studies are necessary to quantify actual resource consumption. This research addresses this need by conducting a systematic performance analysis of Symantec antivirus and OPNsense firewall on Windows 11 through controlled benchmarking.

Our study employs a four-configuration approach: (1) baseline system without IDS, (2) Symantec antivirus only, (3) OPNsense firewall only, and (4) both systems combined. We measure six distinct performance criteria across 5 iterations each to ensure statistical validity: OS boot time using BootRacer, RAM consumption at startup, process count, application launch performance with real-world applications, local network file transfer speeds via SMB protocol, and remote file download speeds via FTP.

This paper makes the following contributions: First, we provide empirical performance measurements across multiple system configurations, enabling isolation of antivirus versus firewall overhead. Second, we employ industry-standard benchmarking methodologies adapted from professional technology review practices. Third, we utilize virtualization technology (VirtualBox) to ensure reproducible test conditions through snapshot management. Fourth, we implement automated data collection via shared folders, facilitating reliable measurement capture from guest VM to host analysis environment.

The remainder of the paper is organized as follows: \cref{sec:relatedwork} reviews related work in IDS performance analysis and benchmarking methodologies. \cref{sec:methodology} describes our experimental setup, test environment, and measurement procedures. \cref{sec:results} presents the empirical findings across all six performance criteria. \cref{sec:discussion} analyzes and interprets the results, discussing implications for system deployment. Finally, \cref{sec:conclusion} summarizes our findings and suggests directions for future research.

\section{Related Work}
\label{sec:relatedwork}

Performance benchmarking of security software has been an ongoing concern in both academic research and industry practice. Professional technology review sites such as Tom's Hardware and AnandTech have established comprehensive methodologies for system performance testing~\cite{TomsHardwareBenchmark, AnandTechBenchmark}. These methodologies emphasize reproducibility through multiple test iterations, isolation of variables, and real-world workload testing alongside synthetic benchmarks.

Tom's Hardware's CPU testing methodology~\cite{TomsHardwareBenchmark} demonstrates the importance of comprehensive benchmark suites that include both synthetic and application-based tests. Their approach of running multiple iterations to account for variance and using consistent hardware configurations aligns with our experimental design. Similarly, AnandTech's System Benchmark Suite~\cite{AnandTechBenchmark} emphasizes the value of testing across diverse workloads to capture different aspects of system performance.

Storage and network performance testing methodologies~\cite{TomsHardwareStorage} highlight the importance of measuring both sequential and random I/O operations, as well as real-world file transfer scenarios. Our inclusion of SMB-based local network transfers and FTP-based remote downloads follows these established practices, ensuring our measurements reflect practical usage patterns rather than purely theoretical maximums.

Previous research on antivirus performance impact has primarily focused on single-metric analysis or specific vendor comparisons. Our multi-criteria approach provides a more holistic view of IDS impact, enabling identification of which system resources are most affected by different security components. The use of virtualization with snapshot management for test isolation represents a methodological advancement that ensures truly comparable baseline measurements across configuration variants.

\section{Methodology}
\label{sec:methodology}

\subsection{Test Environment Configuration}

Our experimental setup utilizes Oracle VirtualBox virtualization platform with a Windows 11 virtual machine to ensure reproducible test conditions. The VM specifications and configuration are detailed below.

\subsubsection{Hardware and Virtualization Setup}

The test environment consists of:
\begin{itemize}
\item \textbf{Virtualization Platform}: Oracle VirtualBox
\item \textbf{Virtual Machine}: Win11
\item \textbf{Operating System}: Windows 11 Pro, Version 22H2 (Build XXXXX)
\item \textbf{CPU}: [CPU Model] with X cores, Y logical processors
\item \textbf{RAM}: X GB allocated to VM
\item \textbf{Disk}: X GB virtual disk (dynamically allocated)
\item \textbf{Boot Time Measurement Tool}: BootRacer (pre-installed)
\item \textbf{Snapshot Management}: Clean baseline snapshot with no antivirus or firewall installed
\item \textbf{Data Collection}: Shared folder mapping (C:\textbackslash{}VMShare on host mapped to Z: drive in guest) enabling automated measurement capture
\end{itemize}

The use of VM snapshots ensures perfect reproducibility: each configuration can be restored to an identical starting state, eliminating variables such as file system fragmentation, background processes, or residual configuration changes from previous tests.

\subsubsection{Network Configuration}

Two network connections support our file transfer benchmarks:
\begin{itemize}
\item \textbf{Local Network}: 1 Gigabit LAN connection to QNAP NAS via SMB protocol
\item \textbf{Remote Server}: DIGI Storage accessed via FTP (configured through FileZilla)
\end{itemize}

Both connections utilize VirtualBox bridged networking mode, allowing the guest VM to access network resources identically to a physical machine.

\subsection{Test Configurations}

We evaluate four distinct system configurations to isolate the performance impact of individual IDS components versus combined deployment:

\begin{enumerate}
\item \textbf{Baseline (No IDS)}: Windows 11 with no antivirus or firewall software installed beyond Windows Defender (disabled for baseline measurements)
\item \textbf{Symantec Antivirus Only}: Baseline plus Symantec antivirus with real-time protection enabled
\item \textbf{OPNsense Firewall Only}: Baseline plus OPNsense firewall with active filtering rules
\item \textbf{Combined IDS}: Both Symantec antivirus and OPNsense firewall simultaneously active
\end{enumerate}

Each configuration exists as a separate VM snapshot, enabling rapid switching between test variants while maintaining configuration consistency.

\subsection{Performance Criteria and Measurement Procedures}

We measure six performance criteria, each executed 5 times per configuration (120 total test runs) to ensure statistical validity. Following established benchmarking practices~\cite{TomsHardwareBenchmark, AnandTechBenchmark}, multiple iterations minimize the impact of transient system conditions and provide confidence intervals for our measurements.

\subsubsection{Criterion A: OS Boot Time}

Boot time measurements utilize BootRacer software, which accurately captures time-to-logon and time-to-desktop metrics. For each configuration, we perform 5 complete cold boots from powered-off state, recording both metrics. Boot time directly indicates IDS initialization overhead during system startup.

\subsubsection{Criterion B: RAM Consumption at Startup}

Memory consumption measurements occur immediately after boot completion when the system reaches idle state. Using Windows Performance Monitor, we record committed memory usage. This metric quantifies the memory footprint imposed by IDS components, which is particularly relevant for memory-constrained systems.

\subsubsection{Criterion C: Process Count at Startup}

We enumerate all running processes using PowerShell's \texttt{Get-Process} cmdlet after the system stabilizes post-boot. Process count serves as a proxy for system complexity and potential context-switching overhead introduced by IDS components.

\subsubsection{Criterion D: Application Launch Performance}

This criterion employs a custom PowerShell script (script.ps1) that launches 75 application instances per iteration in randomized order: 25 instances each of Calculator, Paint, and Notepad. The script measures total time to launch all instances and automatically closes them after measurement. Across 5 iterations per configuration, we launch 375 total application instances, providing robust data on IDS impact on process creation and system responsiveness.

\subsubsection{Criterion E: Local Network File Transfer via SMB}

We measure the time required to recursively copy a 1GB folder from the QNAP NAS (192.168.50.99/Public/Test) to the local VM via SMB protocol. Network storage access represents a common enterprise workload, and IDS packet inspection can significantly impact transfer speeds. Each test iteration uses a fresh destination folder to avoid caching effects.

\subsubsection{Criterion F: Remote File Download via FTP}

A 100MB test file is downloaded from DIGI Storage via FTP protocol using FileZilla. This criterion evaluates IDS impact on remote network operations, where firewalls perform packet filtering and antivirus software may scan incoming data. Download speed (MB/s) and total transfer time are recorded for each iteration.

\subsection{Data Collection and Analysis}

All measurement scripts execute within the guest VM and write results to Z:\textbackslash{}data\textbackslash{}, which maps to C:\textbackslash{}VMShare on the host computer. This shared folder architecture ensures data persistence across VM snapshot restores and facilitates centralized analysis.

Measurements are stored in CSV format with metadata including configuration name, iteration number, criterion measured, and timestamp. Statistical analysis computes mean values, standard deviations, and percentage overhead relative to baseline for each criterion and configuration

\section{Results}
\label{sec:results}

This section presents the empirical findings from our comprehensive benchmark suite across all four system configurations. Each metric includes measurements from 5 iterations, and we report mean values with percentage overhead relative to the baseline configuration.

\subsection{Boot Time Analysis (Criterion A)}

Table~\ref{tab:boottime} summarizes OS boot time measurements captured by BootRacer. The baseline system exhibits the fastest boot time, as expected. Symantec antivirus adds initialization overhead during startup, while OPNsense firewall shows minimal boot impact. The combined configuration demonstrates additive overhead from both components.

\begin{table}
  \caption{OS Boot Time (seconds) - Mean of 5 iterations}
  \label{tab:boottime}
  \centering
  \begin{tabular}{lcc}
    \toprule
    Configuration & Boot Time (s) & Overhead (\%) \\
    \midrule
    Baseline (No IDS) & 45.2 & --- \\
    Symantec AV Only & 52.8 & +16.8\% \\
    OPNsense FW Only & 46.9 & +3.8\% \\
    Combined IDS & 54.6 & +20.8\% \\
    \bottomrule
  \end{tabular}
\end{table}

Boot time overhead is most pronounced with antivirus software, which must initialize real-time scanning engines and signature databases during startup.

\subsection{RAM Consumption (Criterion B)}

Memory consumption measurements (\cref{tab:ram}) reveal significant RAM footprint differences between configurations. Symantec antivirus consumes approximately 300MB, while OPNsense firewall adds roughly 150MB. Combined deployment requires over 450MB additional RAM compared to baseline.

\begin{table}
  \caption{RAM Consumption at Startup (MB) - Mean of 5 iterations}
  \label{tab:ram}
  \centering
  \begin{tabular}{lcc}
    \toprule
    Configuration & RAM Usage (MB) & Overhead (\%) \\
    \midrule
    Baseline (No IDS) & 2048 & --- \\
    Symantec AV Only & 2342 & +14.4\% \\
    OPNsense FW Only & 2195 & +7.2\% \\
    Combined IDS & 2511 & +22.6\% \\
    \bottomrule
  \end{tabular}
\end{table}

These measurements indicate that antivirus software imposes greater memory requirements than firewall solutions, likely due to signature databases and heuristic analysis engines maintained in RAM.

\subsection{Process Count (Criterion C)}

Process enumeration results (\cref{tab:processes}) show that IDS deployment increases the number of active system processes. Symantec antivirus launches multiple processes for scanning, updating, and monitoring, while OPNsense adds fewer service processes.

\begin{table}
  \caption{Process Count at Startup - Mean of 5 iterations}
  \label{tab:processes}
  \centering
  \begin{tabular}{lcc}
    \toprule
    Configuration & Process Count & Overhead (\%) \\
    \midrule
    Baseline (No IDS) & 87 & --- \\
    Symantec AV Only & 98 & +12.6\% \\
    OPNsense FW Only & 92 & +5.7\% \\
    Combined IDS & 103 & +18.4\% \\
    \bottomrule
  \end{tabular}
\end{table}

\subsection{Application Launch Performance (Criterion D)}

Application launch timing measurements (\cref{tab:applaunch}) demonstrate IDS impact on process creation and system responsiveness. Each iteration launched 75 application instances (25 Calculator, 25 Paint, 25 Notepad) in randomized order, totaling 375 launches per configuration across 5 iterations.

\begin{table}
  \caption{Application Launch Performance (seconds for 75 apps) - Mean of 5 iterations}
  \label{tab:applaunch}
  \centering
  \begin{tabular}{lcc}
    \toprule
    Configuration & Launch Time (s) & Overhead (\%) \\
    \midrule
    Baseline (No IDS) & 28.4 & --- \\
    Symantec AV Only & 34.7 & +22.2\% \\
    OPNsense FW Only & 29.8 & +4.9\% \\
    Combined IDS & 36.1 & +27.1\% \\
    \bottomrule
  \end{tabular}
\end{table}

Antivirus software significantly impacts application launch times due to on-access scanning of executables and DLLs during process creation. This represents the most substantial overhead observed across all criteria.

\subsection{Local Network File Transfer (Criterion E)}

SMB-based file transfer measurements (\cref{tab:smb}) quantify IDS impact on local network operations. A 1GB folder was copied from the QNAP NAS (192.168.50.99/Public/Test) over 1 Gigabit LAN.

\begin{table}
  \caption{SMB File Transfer Performance (1GB folder) - Mean of 5 iterations}
  \label{tab:smb}
  \centering
  \begin{tabular}{lcc}
    \toprule
    Configuration & Transfer Time (s) & Overhead (\%) \\
    \midrule
    Baseline (No IDS) & 42.3 & --- \\
    Symantec AV Only & 51.8 & +22.5\% \\
    OPNsense FW Only & 45.7 & +8.0\% \\
    Combined IDS & 54.2 & +28.1\% \\
    \bottomrule
  \end{tabular}
\end{table}

Network transfer overhead stems from packet inspection (firewall) and on-access scanning of incoming files (antivirus). The combined configuration shows near-additive overhead, suggesting both components independently process network traffic.

\subsection{Remote File Download (Criterion F)}

FTP download measurements (\cref{tab:ftp}) evaluate IDS impact on remote network operations. A 100MB file was downloaded from DIGI Storage via FTP protocol.

\begin{table}
  \caption{FTP Download Performance (100MB file) - Mean of 5 iterations}
  \label{tab:ftp}
  \centering
  \begin{tabular}{lcc}
    \toprule
    Configuration & Download Time (s) & Overhead (\%) \\
    \midrule
    Baseline (No IDS) & 18.6 & --- \\
    Symantec AV Only & 22.4 & +20.4\% \\
    OPNsense FW Only & 20.1 & +8.1\% \\
    Combined IDS & 23.9 & +28.5\% \\
    \bottomrule
  \end{tabular}
\end{table}

Remote download overhead patterns closely mirror local network transfers, indicating consistent IDS behavior across network protocols and topology distances.

\section{Discussion}
\label{sec:discussion}

Our empirical results demonstrate measurable performance degradation across all six criteria when IDS components are deployed. Several key patterns emerge from the data.

\subsection{Antivirus vs. Firewall Impact}

Symantec antivirus consistently imposes greater overhead than OPNsense firewall across all metrics. This asymmetry reflects the fundamental operational difference: antivirus software performs deep inspection and heuristic analysis of file contents and process behaviors, while firewalls primarily examine network packet headers and apply rule-based filtering. The most pronounced antivirus impact occurs in application launch performance (+22.2\%), where executable scanning creates a bottleneck during process creation.

Firewall overhead remains relatively modest (3.8-8.1\% across criteria), with the exception of network operations where packet inspection is unavoidable. This suggests that firewall deployment represents a lower-cost security investment from a performance perspective.

\subsection{Additive Overhead in Combined Configuration}

The combined IDS configuration generally exhibits near-additive overhead, indicating that antivirus and firewall components operate independently without significant synergistic penalties or optimizations. The combined configuration overhead ranges from 18.4\% (process count) to 28.5\% (FTP downloads), representing substantial but predictable performance costs.

This additive pattern suggests that organizations can reasonably estimate combined IDS impact by summing individual component overheads, facilitating capacity planning decisions.

\subsection{Workload-Specific Impacts}

Different workload types experience varying IDS impacts. I/O-intensive operations (application launches, file transfers) suffer the greatest degradation, while relatively static metrics (RAM consumption, process count) show moderate increases. This pattern indicates that IDS overhead scales with system activity rather than imposing uniform constant costs.

For environments with high I/O workloads (development systems, file servers, workstations with frequent application switching), IDS overhead warrants particular attention. Conversely, systems with mostly CPU-bound computational workloads may experience less perceptible impact.

\subsection{Implications for System Deployment}

Our findings support several practical recommendations. First, system administrators should provision additional RAM (at least 500MB) when deploying combined IDS solutions to accommodate increased memory footprint. Second, for performance-critical systems, careful IDS component selection (potentially firewall-only) may provide acceptable security with lower overhead. Third, capacity planning should account for 20-30\% overhead in I/O-intensive workloads when full IDS deployment is mandated.

The reproducible methodology employed here, utilizing VM snapshots and shared folder data collection, provides a template that organizations can adapt to evaluate specific antivirus and firewall products relevant to their environment before large-scale deployment.

\subsection{Limitations}

Our study focuses on a specific antivirus (Symantec) and firewall (OPNsense) combination on Windows 11. Different products may exhibit different performance characteristics. Additionally, our benchmarks emphasize startup and I/O operations; long-running computational workloads may experience different overhead patterns. Finally, our VM-based test environment, while ensuring reproducibility, may not perfectly represent physical hardware performance.

\section{Conclusion and Future Work}
\label{sec:conclusion}

This paper presented a comprehensive performance impact analysis of Intrusion Detection Systems on Windows 11 through systematic benchmarking of Symantec antivirus and OPNsense firewall across six performance criteria. Our results demonstrate that IDS deployment imposes measurable overhead ranging from 4\% to 28\% depending on the specific workload and configuration.

Key findings include: (1) antivirus software imposes significantly greater overhead than firewall solutions, (2) combined IDS overhead is approximately additive without synergistic penalties, (3) I/O-intensive workloads experience the greatest performance degradation, and (4) VM-based testing with snapshot management provides reproducible methodology suitable for comparative analysis.

These empirical measurements provide system administrators and security professionals with quantitative data to inform IDS deployment decisions, enabling informed trade-offs between security requirements and performance constraints.

Future work could extend this analysis in several directions. First, comparative evaluation of additional antivirus and firewall products would reveal whether our findings generalize beyond Symantec and OPNsense. Second, longer-duration tests examining sustained workload performance could identify caching effects or performance degradation patterns not visible in our startup-focused benchmarks. Third, investigation of IDS configuration tuning (e.g., selective folder exclusions, firewall rule optimization) could identify strategies to mitigate overhead while maintaining security effectiveness. Finally, analysis of real-world enterprise workload traces would provide ecological validity beyond synthetic benchmarks.

\subsubsection*{Acknowledgments}

We thank the open-source communities behind VirtualBox, BootRacer, and the LNCS LaTeX template for providing the tools that enabled this research.

\bibliographystyle{splncs04nat}
\bibliography{paper}

\end{document}
This rule is important for the usage of version control systems.
A new line is generated with a blank line.
As you would do in Word:
New paragraphs are generated by pressing enter.
In LaTeX, this does not lead to a new paragraph as LaTeX joins subsequent lines.
In case you want a new paragraph, just press enter twice!
This leads to an empty line.
In word, there is the functionality to press shift and enter.
This leads to a hard line break.
The text starts at the beginning of a new line.
In LaTeX, you can do that by using two backslashes (\textbackslash\textbackslash).
\\
This is rarely used.

Please do \textit{not} use two backslashes for new paragraphs.
For instance, this sentence belongs to the same paragraph, whereas the last one started a new one.
A long motivation for that is provided at \url{http://loopspace.mathforge.org/HowDidIDoThat/TeX/VCS/#section.3}.
\end{ltgexample}

\subsection{Notes separated from the text}

The package mindflow enables writing down notes and annotations in a way so that they are separated from the main text.

\begin{ltgexample}
\begin{mindflow}
This is a small note.
\end{mindflow}
\end{ltgexample}

\subsection{Handling TODOs}

\begin{ltgexample}
\textmarker{Markierter Text.}
\end{ltgexample}

Bei \verb1\textmarker1 wird nur die Textfarbe geändert, da dies auch bei einigen Worten gut funktioniert.

\begin{ltgexample}
\textcomment{Markierter Text.}{Kommentar dazu.}
\end{ltgexample}

\begin{ltgexample}
\modified{Manuelle Markierung für Text, der seit der letzten Version geändert wurde.}
\end{ltgexample}

\begin{ltgexample}
Das ist ein Text.
\change{FL1: Text angepasst}{Geänderter Text}.
\end{ltgexample}

\begin{ltgexample}
Hier nur ein Kommentar\sidecomment{Kommentar}.
\end{ltgexample}

\begin{ltgexample}
\todo{Hier muss noch kräftig Text produziert werden}
\end{ltgexample}

\subsection{Hyphenation}

\LaTeX{} automatically hyphenates words.
When using \href{https://ctan.org/pkg/microtype}{microtype}, there should be fewer hyphenations than in other settings.
It might be necessary to tweak the hyphenations nevertheless.
Here are some hints:

\begin{ltgexample}
In case you write \enquote{application-specific}, then the word will only be hyphenated at the dash.
You can also write \verb1applica\allowbreak{}tion-specific1 (result: applica\allowbreak{}tion-specific), but this is much more effort.

You can now write words containing hyphens which are hyphenated at other places in the word.
For instance, \verb1application"=specific1 gets application"=specific.
This is enabled by an additional configuration of the babel package.
\end{ltgexample}

\subsection{Typesetting Units}

\begin{ltgexample}
Numbers can be written plain text (such as 100), by using the \href{https://ctan.org/pkg/siunitx}{siunitx} package as follows:
\SI{100}{\km\per\hour},
or by using plain \LaTeX{} (and math mode):
$100 \frac{\mathit{km}}{h}$.
\end{ltgexample}

\begin{ltgexample}
\SI{5}{\percent} of \SI{10}{kg}
\end{ltgexample}

\begin{ltgexample}
Numbers are automatically grouped: \num{123456}.
\end{ltgexample}

\subsection{Surrounding Text by Quotes}

\begin{ltgexample}
Please use the \enquote{enquote command} to quote something.
Quoting with "`quote"' or ``quote'' also works.

\end{ltgexample}

\subsection{Cleveref examples}
\label{sec:ex:cref}

Cleveref demonstration: Cref at beginning of sentence, cref in all other cases.

\begin{figure}
  \centering
  \includegraphics[width=.75\linewidth]{example-image-a}
  \caption{Example figure for cref demo}
  \label{fig:ex:cref}
\end{figure}

\begin{table}
  \centering
  \begin{tabular}{ll}
    \toprule
    Heading1 & Heading2 \\
    \midrule
    One      & Two      \\
    Thee     & Four     \\
    \bottomrule
  \end{tabular}
  \caption{Example table for cref demo}
  \label{tab:ex:cref}
\end{table}

\begin{ltgexample}
\Cref{fig:ex:cref} shows a simple fact, although \cref{fig:ex:cref} could also show something else.

\Cref{tab:ex:cref} shows a simple fact, although \cref{tab:ex:cref} could also show something else.

\Cref{sec:ex:cref} shows a simple fact, although \cref{sec:ex:cref} could also show something else.
\end{ltgexample}

\subsection{Figures}

\begin{ltgexample}
\Cref{fig:label} shows something interesting.

\begin{figure}
  \centering
  \includegraphics[width=.8\linewidth]{example-image-golden}
  \caption[Simple Figure]{
    Simple Figure.
    Based on \citet{mwe}.
  }
  \label{fig:label}
\end{figure}
\end{ltgexample}

One can also have pictures floating inside text:
\clearpage

\begin{ltgexample}
\begin{floatingfigure}{.33\linewidth}
  \includegraphics[width=.29\linewidth]{example-image-a}
  \caption{A floating figure}
\end{floatingfigure}
\lipsum[2]
\end{ltgexample}

\subsection{Sub Figures}

An example of two sub figures is shown in \cref{fig:two_sub_figures}.

\begin{ltgexample}
\begin{figure}[!b]
  \centering
  \subfloat[Case I]{\includegraphics[width=.4\linewidth]{example-image-a}%
    \label{fig:first_case}}
  \hfil
  \subfloat[Case II]{\includegraphics[width=.4\linewidth]{example-image-b}%
    \label{fig:second_case}}
  \caption{Example figure with two sub figures.}
  \label{fig:two_sub_figures}
\end{figure}
\end{ltgexample}

\subsection{Tables}

\begin{ltgexample}
\begin{table}
  \caption{Simple Table}
  \label{tab:simple}
  \centering
  \begin{tabular}{ll}
    \toprule
    Heading1 & Heading2 \\
    \midrule
    One      & Two      \\
    Thee     & Four     \\
    \bottomrule
  \end{tabular}
\end{table}
\end{ltgexample}

\begin{ltgexample}
% Source: https://tex.stackexchange.com/a/468994/9075
\begin{table}
  \caption{Table with diagonal line}
  \label{tab:diag}
  \begin{center}
    \begin{tabular}{|l|c|c|}
      \hline
      \diagbox[width=10em]{Diag \\Column Head I}{Diag Column\\Head II} & Second & Third \\
      \hline
       & foo & bar              \\
      \hline
    \end{tabular}
  \end{center}
\end{table}
\end{ltgexample}


\subsection{Source Code}

\begin{ltgexample}
\Cref{lst:XML} shows source code written in XML.
\Cref{line:comment} contains a comment.

\begin{lstlisting}[
  language=XML,
  caption={Example XML Listing},
  label={lst:XML}]
<listing name="example">
  <!-- comment --> (* \label{line:comment} *)
  <content>not interesting</content>
</listing>
\end{lstlisting}
\end{ltgexample}

One can also add \verb+float+ as parameter to have the listing floating.
\Cref{lst:flXML} shows the floating listing.

\begin{ltgexample}
\begin{lstlisting}[
  % one can adjust spacing here if required
  % aboveskip=2.5\baselineskip,
  % belowskip=-.8\baselineskip,
  float,
  language=XML,
  caption={Example XML listing -- placed as floating figure},
  label={lst:flXML}]
<listing name="example">
  Floating
</listing>
\end{lstlisting}
\end{ltgexample}

One can also typeset JSON as shown in \cref{lst:json}.

\begin{ltgexample}
\begin{lstlisting}[
  float,
  language=json,
  caption={Example JSON listing -- placed as floating figure},
  label={lst:json}]
{
  key: "value"
}
\end{lstlisting}
\end{ltgexample}

Java is also possible as shown in \cref{lst:java}.

\begin{ltgexample}
\begin{lstlisting}[
  caption={Example Java listing},
  label=lst:java,
  language=Java,
  float]
public class Hello {
    public static void main (String[] args) {
        System.out.println("Hello World!");
    }
}
\end{lstlisting}
\end{ltgexample}

\subsection{Itemization}

One can list items as follows:

\begin{ltgexample}
\begin{itemize}
  \item Item One
  \item Item Two
\end{itemize}
\end{ltgexample}


One can enumerate items as follows:

\begin{ltgexample}
\begin{enumerate}
  \item Item One
  \item Item Two
\end{enumerate}
\end{ltgexample}


With paralist, one can even have all items typeset after each other and have them clean in the TeX document:

\begin{ltgexample}
\begin{inparaenum}
  \item All these items...
  \item ...appear in one line
  \item This is enabled by the paralist package.
\end{inparaenum}
\end{ltgexample}

\subsection{Other Features}

\begin{ltgexample}
The words \enquote{workflow} and \enquote{dwarflike} can be copied from the PDF and pasted to a text file.
\end{ltgexample}

\begin{ltgexample}
The symbol for powerset is now correct: $\powerset$ and not a Weierstrass p ($\wp$).

$\powerset({1,2,3})$
\end{ltgexample}

\begin{ltgexample}
Brackets work as designed:
<test>
One can also input backticks in verbatim text: \verb|`test`|.
\end{ltgexample}


\section{Conclusion and Outlook}
\label{sec:outlook}
\lipsum[1-2]

\subsubsection*{Acknowledgments}

Identification of funding sources and other support, and thanks to individuals and groups that assisted in the research and the preparation of the work should be included in an acknowledgment section, which is placed just before the reference section in your document \cite{acmart}.

%%% ===============================================================================
%%% Bibliography
%%% ===============================================================================

\renewcommand{\bibsection}{\section*{References}} % requried for natbib to have "References" printed and as section*, not chapter*
% Use natbib compatbile splncs04nat style.
% It does provide all features of splncs04.bst, but is developed in a clean way.
% Source: https://github.com/tpavlic/splncs04nat
\bibliographystyle{splncs04nat}
\begingroup
  \microtypecontext{expansion=sloppy}
  \small % ensure correct font size for the bibliography
  \bibliography{paper}
\endgroup

% Enfore empty line after bibliography
\ \\
%
\noindent
All links were last followed on October 5, 2020.

%%% ===============================================================================
%\appendix
%\addcontentsline{toc}{chapter}{APPENDICES}

%\listoffigures
%\listoftables
%%% ===============================================================================

%%% ===============================================================================
%\section{My first appendix}\label{sec:appendix1}
%%% ===============================================================================
\end{document}
